% ------------------------------------------------------------------------
%
\begin{frame}{Unification}

  The best way to think about unification is to think of terms as
  trees, as we did page~\pageref{terms_as_trees}. A function is a
  subtree and a constant is a leaf.

  \bigskip

  The unification algorithm is then very simple:
  \begin{itemize}

    \item if the two terms are constants, they match if they are
      verbatim the same;

    \item if one of the terms is a variable, then they match;

    \item if the two terms are functions, they match if and only if
      \begin{itemize}

        \item they are the same (name),

        \item they have the same \textbf{arity} (number of arguments,
         that is, number of subterms, or number of subtrees),

        \item if they have arguments, the \(n\)-th argument of one
          term matches the \(n\)-th argument of the other.

  \end{itemize}

\end{itemize}

\end{frame}

% ------------------------------------------------------------------------
%
\begin{frame}{Unification / Example \#1}

  For example, let's unify the terms \(\fun{f}(\fun{g}(x,x),
  \fun{h}(y,y))\) and \(\fun{f}(\fun{g}(z), z)\).
  \begin{center}
    \includegraphics{tree-1-1}
  \end{center}
  They do not unify because the subterm \(\fun{g}(x,x)\) does not
  match \(g(z)\): the arity of \fun{g} is different.

  \bigskip

  Note how the terms are non-linear, even though we do not allow them
  in our systems. Indeed, unification is all about equalities, which
  non\hyp{}linearity is about. So let's allow it for exercises.

\end{frame}

% ------------------------------------------------------------------------
%
\begin{frame}{Unification / Example \#2}

  Consider now the terms \(\fun{f}(\fun{g}(x, x), \fun{h}(y, y))\) and
  \(\fun{f}(\fun{g}(\fun{h}(w, \fun{a}()), z), z))\):
  \begin{center}
    \includegraphics{tree-1-2}
  \end{center}

\end{frame}

% ------------------------------------------------------------------------
%
\begin{frame}{Unification / Example \#2}

  The steps are
\begin{align*}
  x &= \fun{h}(w, \fun{a}()) & x &= \fun{h}(w, \fun{a}())
  & \fun{h}(w, \fun{a}()) &= \fun{h}(y, y) & w &= y\\
x &= z      & x &= \fun{h}(y, y) & x      &= \fun{h}(y, y) & x &= \fun{h}(y, y)\\
z &= \fun{h}(y, y) & z &= x      & z      &= x      & y &= \fun{a}()\\
  &                &   &         &        &         & z &= x
\end{align*}
  \noindent Finally we obtain the \textbf{substitution}
\begin{align*}
w &= \fun{a}()\\
x &= \fun{h}(\fun{a}(), \fun{a}())\\
y &= \fun{a}()\\
z &= \fun{h}(\fun{a}(), \fun{a}())
\end{align*}

\end{frame}

% ------------------------------------------------------------------------
%
\begin{frame}{Unification / Example \#2}

  This answer means that if we take the two trees of the query and
  replace the variables according to the substitution, we get the
  exact same tree. Indeed, we get
  \begin{center}
    \includegraphics{tree-1-2u}
  \end{center}

\end{frame}

%% % ------------------------------------------------------------------------
%% %
%% \begin{frame}{Unification / Example \#3}

%%     Consider now the terms \(\fun{f}(\fun{g}(x,x), \fun{h}(\_,\_))\)
%%     and \(\fun{f}(\fun{g}(\fun{h}(w, a()),z), z)\), where the
%%     underscores denote \textbf{fresh variables} (that is, they are
%%     unknown but different from any other variable in the term):
%%     \begin{center}
%%       \includegraphics{tree-1-3}
%%     \end{center}

%% \end{frame}

%% % ------------------------------------------------------------------------
%% %
%% \begin{frame}{Unification / Example \#3}

%%   The steps are now
%%   \begin{align*}
%%     x &= \fun{h}(w,\fun{a}()) & x &= \fun{h}(w, \fun{a}())
%%       & x &= \fun{h}(w, \fun{a}()) & x &= \fun{h}(\alpha, \fun{a}())\\
%%     x &= z & x &= z & z &= x & z &= \fun{h}(\alpha, \fun{a}())\\
%%     \fun{h}(\alpha, \beta) &= z &
%%     \fun{h}(\alpha, \beta) &= \fun{h}(w, \fun{a}())
%%     & w &= \alpha & w &= \alpha\\
%%     & & & & & & \beta &= \fun{a}()
%%   \end{align*}

%%   Finally
%%   \begin{align*}
%%     \beta &= \alpha \\
%%     w &= \alpha\\
%%     x &= \fun{h}(\alpha, \fun{a}())\\
%%     z &= \fun{h}(\alpha, \fun{a}())
%%   \end{align*}

%% \end{frame}

%% % ------------------------------------------------------------------------
%% %
%% \begin{frame}{Unification / Example \#3}

%%   This means that if we replace the variables in the two initial
%%   trees, we get the same tree. Indeed, we get
%%   \begin{center}
%%     \includegraphics{tree-1-3u}
%%   \end{center}

%% \end{frame}
